\chapter{Annexes}

\begin{lstlisting}[language=Python, caption={Extraction OCR – conversion d’une page PDF en image puis reconnaissance par Tesseract.}, label={code:extrait-ocr}]
def extraire_avec_ocr_neutre(self, pdf_path):
    """OCR sans interprétation linguistique"""
    text_parts = []
    try:
        doc = fitz.open(pdf_path)
        for page_num in range(len(doc)):
            page = doc.load_page(page_num)
            # Vérifier d'abord si la page a du texte vectoriel
            text_vectoriel = page.get_text()
            if text_vectoriel and len(text_vectoriel.strip()) > self.min_text_length:
                continue
            
            # Sinon, faire OCR
            try:
                pix = page.get_pixmap(matrix=fitz.Matrix(300/72, 300/72))
                img_data = pix.tobytes("ppm")
                
                from io import BytesIO
                img = Image.open(BytesIO(img_data))
                
                # Convertir en niveaux de gris pour meilleur OCR
                if img.mode != 'L':
                    img = img.convert('L')
                
                # OCR avec configuration neutre
                text = pytesseract.image_to_string(
                    img, 
                    config=self.ocr_config
                )
                if text.strip():
                    text_parts.append(f"\n{'='*40} Page {page_num + 1} (OCR) {'='*40}\n")
                    text_parts.append(text)
            except Exception as ocr_error:
                print(f"    OCR page {page_num + 1} échoué: {ocr_error}")
                continue
        doc.close()
        return "\n".join(text_parts)
    except Exception as e:
        print(f"  Erreur OCR neutre: {e}")
        return ""
\end{lstlisting}

